\section{\bench 与评测协议}
\label{app:benchmark}
\suppressfloats[t]

\subsection{任务、角色与成功条件}
\label{app:tasks}
我们在 MuJoCo 与 robosuite \cite{mujoco,robosuite} 上构建了 \bench。它包含五个任务族共 14 个任务（图~\mainref{fig:bench_overview}）。表~\ref{tab:tasks} 列出它们的成功条件，表~\ref{tab:task_sampling} 列出其场景范围。我们将公开这些任务、已接受的场景、成功判定与示范。

\begin{table*}[!t]
\centering\fontsize{9}{11}\selectfont
\setlength{\tabcolsep}{4pt}
\renewcommand{\arraystretch}{1.13}
\caption{Native success conditions in the reported fourteen-task suite.
Thresholds come from the evaluator source and task configurations bound to the
accepted FUTURE0, BASE2, and HISTORY1 packages. A task must pass its preceding
stages and avoid terminal failures. All tasks run at 20\,Hz with a 2\,ms physics
step; durations below refer to physical time. COM denotes center of mass.}
\label{tab:tasks}
\begin{tabularx}{\textwidth}{@{}p{0.19\textwidth}Xp{0.07\textwidth}@{}}
\toprule
Task & Ordered success conditions & Limit \\
\midrule
grasp;\newline t1\_belt, t1\_xyz,\newline t1\_so3, t1\_occ
& Approach within 0.10\,m; bilateral fingerpad contact for 0.05\,s; collision-geometry
bottom at least 0.05\,m above support; retain contact and clearance for 0.50\,s.
T1 adds conveyor drive, translation, rotation, or a physical occluder.
SO3 has no additional success-angle threshold; unsafe occluder contact terminates OCC.
& 160 steps\newline 8\,s \\
t2\_trans
& Approach within 0.10\,m. Entire object collision geometry inside the receptacle:
planar clearance and floor gap both $\geq-0.002$\,m, top no higher than the opening.
Release requires an opening command $<-0.5$ and no fingerpad contact.
Relative linear/angular speeds $\leq0.08$\,m/s and $0.8$\,rad/s for 0.10\,s,
followed by a 0.50\,s stable dwell.
& 200\newline 10\,s \\
t2\_se3
& Approach within 0.05\,m and 0.12\,rad. Full T-shaped containment:
planar clearance $\geq0.002$\,m, floor gap in $[-0.002,0.002]$\,m,
tilt $\leq0.05$\,rad, top below opening.
After release, relative speeds $\leq0.02$\,m/s and $0.04$\,rad/s
for 0.10\,s, then a 0.50\,s dwell. The 0.12\,rad gate applies to approach.
& 200\newline 10\,s \\
t3\_phase
& Start with the native pre-grasp. Phase alignment: distance $\leq0.090$\,m and
loop-direction error $\leq0.20$\,rad. Enter within 0.030\,m at progress $\leq0.15$;
reach historical maximum progress $\geq0.60$, then $\geq0.85$ and endpoint
distance $\leq0.055$\,m.
& 420\newline 21\,s \\
t3\_full
& Acquire the wand with bilateral contact for 0.05\,s, then complete the phase
stages above. In a later sample, retain current progress $\geq0.85$ and move
at least 0.045\,m from both path and endpoint. Finally open ($<-0.5$)
with no fingerpad contact. Releasing before exit fails.
& 640\newline 32\,s \\
t4\_sphere
& Approach within 0.10\,m; first qualifying contact after ramp exit; any tray contact before success
fails. Bilateral hold for nominal 25\,ms (12 physics samples, 24\,ms).
Raise COM by 0.05\,m from its height when the bilateral-hold criterion is first met, and retain
contact and height for 0.50\,s.
& 120\newline 6\,s \\
t4\_sphere\_tabletop;\newline t4\_slider
& Approach within 0.10\,m; contact after ramp exit, while supported, in the capture region, and moving
at planar speed $\geq0.08$\,m/s. Bilateral hold for 0.05\,s, COM lift
$\geq0.05$\,m from the first attainment of that hold criterion, then contact and height for 0.50\,s.
& 120\newline 6\,s \\
t4\_can
& Same supported-capture sequence. Before first fingerpad contact, a can-axis
absolute dot product with vertical $>0.45$ fails as unstable tipping.
The configuration does not require a transverse closing axis.
& 120\newline 6\,s \\
t5\_replan
& Exactly one selected-side impact group completes before robot contact.
At 30\,ms after the physical impact-end timestamp, branch-signed lateral speed $\geq0.035$\,m/s and yaw rate
$\geq0.35$\,rad/s. Approach within 0.10\,m; bilateral hold with grasp-axis
error $\leq18^\circ$ for 0.05\,s; COM lift $\geq0.05$\,m from the first
attainment of that hold criterion and stable hold
for 0.50\,s. Impact-group lag is 12\,ms and maximum contact is 0.10\,s.
& 120\newline 6\,s \\
\bottomrule
\end{tabularx}
\end{table*}


\paragraph{角色。}
\label{app:roles}
表~\ref{tab:roles} 把每个任务中的物体分配到第~\mainref{sec:formulation}~节的目标与环境两种角色。移动的容器是目标，因为它的位姿决定能否放入。穿越任务中的导线同时扮演两种角色：入口处的位姿是目标，因为棒必须在 30\,mm 内与之对齐；沿环的几何是环境，因为任何接触都算失败。传送带、斜坡与撞击器是环境中的驱动物。最后一列列出每个任务涉及的式~\mainref{eq:budget} 中的误差项（附录~\ref{app:analysis}）。图~\mainref{fig:analysis} 的两个任务组由这张表导出，并在消融之前确定：旋转敏感组包含误差项含 $r\theta$ 的任务（\texttt{grasp}、\texttt{t1\_so3}、\texttt{t2\_se3}、\texttt{t3\_phase}、\texttt{t3\_full}、\texttt{t4\_can}、\texttt{t5\_replan}）；环境活跃组包含有障碍间隙或驱动物的任务（\texttt{t1\_belt}、\texttt{t1\_occ}、\texttt{t3\_phase}、\texttt{t3\_full}，以及四个捕获任务与 \texttt{t5\_replan}）。

% Protocol table: interaction roles, tolerances, speeds, critical leads, and exercised budget terms per Dynamic-14 task (from the task definitions and sampling ranges; bracketed values are computed from the task geometry before testing).
\begin{table*}[t]
\centering\footnotesize
\caption{\tablelang{\textbf{Roles and tolerances of the \bench tasks.} The target is the body whose pose error
enters a contact or containment tolerance $\varepsilon$. The context comprises obstacles, avoided with clearance $m$, and drivers that move the target.
Terms: error terms that a task exercises (Eq.~\mainref{eq:budget} and Appendix~\ref{app:analysis}). $v$: moving target. $r\theta$: orientation enters the
contact. $e_{\rm dyn}$: non-constant motion (harmonic, gravity, friction, impact). $m$: a context body must be avoided. Occ.:
history gaps. Delay: command delay built into the task. Groups of Fig.~\mainref{fig:analysis}: R, rotation-sensitive; C, context-active.
$^\dagger$Nominal clearance between the wire and the wand's loop for a centered, perpendicular crossing.}{\textbf{\bench 各任务的角色与容差。}目标是位姿误差进入接触或容纳容差 $\varepsilon$ 的物体；环境包括须以间隙 $m$ 避开的障碍物，以及推动目标运动的驱动物。
误差项：任务涉及的误差分解项（式~\mainref{eq:budget} 与附录~\ref{app:analysis}）。$v$：目标运动。$r\theta$：朝向影响接触。$e_{\rm dyn}$：非匀速运动（谐波、重力、摩擦、撞击）。
$m$：须避开某个环境物体。occ.：历史缺失。delay：任务内置的指令延迟。图~\mainref{fig:analysis} 的任务组：R 为旋转敏感组，C 为环境活跃组。
$^\dagger$导线居中且垂直穿过杆端环时两者之间的名义间隙。}}
\label{tab:roles}
\setlength{\tabcolsep}{3.2pt}
\begin{tabularx}{\textwidth}{@{}lll>{\raggedright\arraybackslash}Xccll@{}}
\toprule
\tablelang{Task}{任务} & \tablelang{Code}{代号} & \tablelang{Target}{目标} & \tablelang{Context}{环境} & $\varepsilon$ (mm) & $m$ (mm) & \tablelang{Terms}{误差项} & \tablelang{Group}{任务组} \\
\midrule
\tablelang{Harmonic grasp}{谐波抓取} & grasp & \tablelang{block}{物块} & \tablelang{table}{桌面} & \fillin{X} & -- & $v$, $r\theta$, $e_{\rm dyn}$ & R \\
\tablelang{Conveyor grasp}{传送带抓取} & t1\_belt & \tablelang{part}{零件} & \tablelang{belt}{传送带} & \fillin{X} & -- & $v$ & C \\
\tablelang{Translating grasp}{平移抓取} & t1\_xyz & \tablelang{part}{零件} & \tablelang{table}{桌面} & \fillin{X} & -- & $v$, $e_{\rm dyn}$ & -- \\
\tablelang{Rotating grasp}{旋转抓取} & t1\_so3 & \tablelang{part ($\omega\le1.2$\,rad/s)}{零件（$\omega\le1.2$\,rad/s）} & \tablelang{table}{桌面} & \fillin{X} & -- & $v$, $r\theta$, $e_{\rm dyn}$ & R \\
\tablelang{Occluded grasp}{遮挡抓取} & t1\_occ & \tablelang{part}{零件} & \tablelang{occluder canopy}{遮挡罩} & \fillin{X} & \fillin{X} & $v$, $m$, occ. & C \\
\tablelang{Moving receptacle}{移动容器} & t2\_trans & \tablelang{receptacle}{容器} & \tablelang{table}{桌面} & \fillin{X} & -- & $v$ & -- \\
\tablelang{Rotating receptacle}{旋转容器} & t2\_se3 & \tablelang{T-shaped receptacle}{T 形容器} & \tablelang{table}{桌面} & \fillin{X} & -- & $v$, $r\theta$ & R \\
\tablelang{Wire passage}{导线穿越} & t3\_phase & \tablelang{wire loop (entry)}{导线环（入口）} & \tablelang{rotating wire}{旋转导线} & 30 & 26.4$^\dagger$ & $v$, $r\theta$, $m$ & R, C \\
\tablelang{Wire passage, full}{导线穿越（完整）} & t3\_full & \tablelang{wire loop (entry)}{导线环（入口）} & \tablelang{rotating wire}{旋转导线} & 30 & 26.4$^\dagger$ & $v$, $r\theta$, $m$ & R, C \\
\tablelang{Ball capture}{小球捕获} & t4\_sphere & \tablelang{ball}{小球} & \tablelang{tray, ramp}{托盘、斜坡} & \fillin{X} & \fillin{X} & $v$, $e_{\rm dyn}$, $m$ & C \\
\tablelang{Ball capture, table}{小球捕获（桌面）} & t4\_sphere\_tabletop & \tablelang{ball}{小球} & \tablelang{table, ramp}{桌面、斜坡} & \fillin{X} & -- & $v$, $e_{\rm dyn}$ & C \\
\tablelang{Slider capture}{滑块捕获} & t4\_slider & \tablelang{slider}{滑块} & \tablelang{table, ramp}{桌面、斜坡} & \fillin{X} & -- & $v$, $e_{\rm dyn}$ & C \\
\tablelang{Can capture}{罐体捕获} & t4\_can & \tablelang{can}{罐体} & \tablelang{table, ramp}{桌面、斜坡} & \fillin{X} & -- & $v$, $r\theta$, $e_{\rm dyn}$ & R, C \\
\tablelang{Post-impact grasp}{撞击后抓取} & t5\_replan & \tablelang{part}{零件} & \tablelang{striker}{撞击器} & \fillin{X} & \fillin{X} & $v$, $r\theta$, $e_{\rm dyn}$, $m$, delay & R, C \\
\bottomrule
\end{tabularx}
\end{table*}


\paragraph{终止事件。}我们在每个 2\,ms 的物理步上，按表~\ref{tab:tasks} 的先后要求判定成功；连续条件一旦中断，其计数器清零。终止事件分为五类。\emph{无效}事件是仿真器状态损坏。\emph{边界}事件是离开工作空间或越过下游、掉落，或目标时间窗过期。\emph{安全}事件是不安全接触、过大的边缘冲量、碰到导线、释放冲量、机器人过早接触、无支撑地接触托盘或被撞击器卡住。\emph{任务违规}是在目标之外释放、偏离路径或倒退、过早释放、偏转而未捕获、倾倒、撞击缺失或多余，或出射运动改变不足。最后一类是\emph{超时}。每个任务族在每个物理步按固定顺序判定事件，并保留第一个终止事件。

\subsection{场景、示范与数据划分}
\label{app:splits}
\begin{table*}[!t]
\centering\fontsize{9}{11}\selectfont
\setlength{\tabcolsep}{4pt}
\renewcommand{\arraystretch}{1.12}
\caption{Historical validation proposal ranges. Length, linear velocity, angle,
and angular velocity use m, m/s, rad, and rad/s unless marked otherwise.
Positions are in world coordinates except the stated ramp-local offsets.
Continuous intervals are uniform proposals; discrete choices are equiprobable
before feasibility checks. Evaluation replays the stored accepted scenarios,
whose distribution need not remain uniform. Harmonic centers are not initial positions.}
\label{tab:task_sampling}
\begin{tabularx}{\textwidth}{@{}p{0.18\textwidth}p{0.34\textwidth}X@{}}
\toprule
Task & Initial-state proposal & Motion or event parameters \\
\midrule
grasp & Harmonic center $c_x\in[-.08,.08]$, $c_y\in[-.26,-.19]$, $z=.822$;
yaw $[-1.1,1.1]$.
& $v_y\in[.025,.125]$, $A_x\in[.015,.035]$, $f_x\in[.3,.7]$\,Hz;
$\omega_z\in[-1.2,1.2]$. \\
t1\_belt & Target $x\in[-.08,.08]$, $y\in[-.26,-.18]$, $z=.822$;
identity orientation, zero initial twist.
& Belt $v_y\in[.04,.10]$. \\
t1\_occ & Same target proposal as Belt; canopy at $(.322,-.29,1.20)$.
& Same belt range; clearance parameter .20; physical occlusion. \\
t1\_xyz & Harmonic center $c_x\in[-.025,.025]$, $c_y\in[-.24,-.20]$, $z=.818$;
yaw $[-.35,.35]$.
& $v_y\in[.04,.09]$, $A_x\in[0,.03]$, $f_x\in[.25,.65]$\,Hz; $\omega_z=0$. \\
t1\_so3 & Same as XYZ.
& Same translation; $\omega_z\in[-1.2,-.4]\cup[.4,1.2]$, equal sign probabilities. \\
t2\_trans & Goal $x\in[-.13,-.08]$, $y\in[-.18,-.12]$, $z=.825$; fixed orientation.
& Goal $v_y\in[.03,.06]$; no rotation. \\
t2\_se3 & Same goal $x,y$ range, $z=.804$; yaw $[-.3,.3]$.
& Goal $v_y\in[.03,.06]$;
$\omega_z\in[-.25,-.08]\cup[.08,.25]$, equal sign probabilities. \\
t3\_phase & Wand $x,y\in[-.03,.03]$, $z=.870$, yaw $[0,2\pi]$;
wire phase $[-.25,.35]$.
& Wire center $(.10,-.02,.94)$; world-X rotation
$\omega_x=-\{2.5,3,\ldots,7.5\}^{\circ}$/s (11 equal-probability values). \\
t3\_full & Wand at $(-.05,-.02,1.04)$ with identity orientation;
same wire phase range.
& Same wire rotation; cradle and release requirement enabled. \\
t4\_sphere & Ramp-local offsets: lateral $[-.015,.015]$, longitudinal/normal
$(-.145,.030)$; lip $(.06,-.02,1.10)$.
& Gravity-driven from rest; gate release at tick 0; slope $[5,15]^\circ$;
target/ramp sliding friction $[.7,.9]$; unsupported capture. \\
t4\_sphere\_tabletop & Lateral $[-.015,.015]$, longitudinal/normal
$(-.124,.034)$; lip $(0,-.06,.808)$.
& Same slope and release; target/ramp friction $[.7,.9]$, table friction .35. \\
t4\_slider & Lateral $[-.015,.015]$, longitudinal/normal $(-.129,.024)$;
supported ramp.
& Same slope and release; target/ramp friction $[.025,.045]$,
table friction $[.012,.018]$. \\
t4\_can & Lateral $[-.010,.010]$, ramp-axis yaw $[-4,4]^\circ$;
longitudinal position follows gate geometry and oriented can extent.
& Same slope and release; target/ramp friction $[.25,.45]$,
table friction $[.09,.15]$. \\
t5\_replan & $x\in[-.005,.005]$, $y=-.29$, $z=.812$;
yaw $[-.08,.08]$; zero initial twist.
& Belt $v_y\in[.12,.14]$; equal-probability left/right striker triggered at
$y=-.13$; stroke .098, pulse .35\,s, actuator scale $[1,1.1]$,
force cap 4\,N; command delay 50\,ms. \\
\bottomrule
\end{tabularx}
\end{table*}

我们对连续范围均匀采样，对离散选项等概率采样，然后进行物理可行性检查。评测回放存储下来的已接受场景。简谐平移遵循 $\vect p(t)=\vect c+(0,v_y,0)t+(A_x\sin(2\pi f_xt+\phi),0,0)$，其中 $\phi\sim\mathcal U[-\pi,\pi]$。捕获任务的目标从静止开始，在 $[5,15]^\circ$ 的斜坡上受重力与摩擦作用运动。只有当目标在开始时至少可见八帧、被完全遮挡 0.5--3.0\,s、之后再次至少可见三帧时，我们才接受一个遮挡场景。

训练池包含 11{,}200 条示范，每个任务 800 条，由使用仿真特权状态的专家生成；学习得到的策略只接收第~\mainref{sec:method}~节的输入。我们以 20\,Hz 记录示范，包括相机图像、深度、机器人状态、下发的指令，以及目标位姿与场景几何的仿真真值，并由此计算未来标签、其有效掩码与指令标签。

我们在 1{,}400 个开发场景（每个任务 100 个）上选择配置。全部 \bench 结果使用另外 1{,}400 个留出场景，它们以相同的范围生成，使用不出现在任何训练、验证或开发回合中的新源种子；按种子递增顺序的前 100 个可行场景构成每个任务的集合。

\subsection{执行协议与延迟测量}
\label{app:sweeps}
\paragraph{执行器。}所有方法都在第~\mainref{sec:inference}~节的执行器下以 50\,ms 控制时钟运行，查询只在控制周期边界开始。第 $t$ 次查询的总延迟 $L_t$ 是查询期间物理推进的时间，即该协议下的推理时间加上任何附加延迟。其动作块的第 $j$ 条指令在式~\mainref{eq:apptime} 的 $\tapp_{t,j}$ 执行，撞击后抓取的 $d_{\rm task}=50$\,ms，其余任务为零。其 $j<k_t$ 的指令（$k_t=\lceil L_t/\Delta_a\rceil$）已经过期并被丢弃，因此第 $t$ 次查询执行第 $j=k_t,\ldots,\min(\cfgChunk,k_t+k_{t+1})-1$ 条指令，每个控制周期一条。当 $k_t+k_{t+1}>\cfgChunk$ 时，例如总延迟恒定且超过 200\,ms，动作块会在下一个到达之前耗尽；若 $k_t\ge\cfgChunk$，其全部指令都已过期。在回合的第一个动作块到达之前，以及动作块耗尽时，控制器保持当前的手臂目标与夹爪指令，物理照常推进，不重复执行任何指令。我们报告的出动作时间是取整到 $k_t\Delta_a$ 之前的实测延迟。

\paragraph{实时协议。}除非另有说明，推理期间物理按该方法的原生延迟推进：每次查询从该方法在 \cfgGPU{} 上以 batch 1 实测的延迟分布中抽取。ACT/DP 原生延迟组在 \cfgGPU{} 上运行，ACT 暂停组、附加延迟组及计划中的其余评测使用 \cfgEvalGPU{}；rollout 本身的计算时间不计入回合。表~\ref{tab:supp_protocols} 报告三种替代协议。在\emph{暂停}协议下，执行器不变，仿真等待推理完成；附加延迟仍推进物理。没有附加延迟时 $L_t=0$，第 0 条指令在采集时刻执行。在\emph{等延迟}协议下，无论自身速度如何，每个方法都被计入 \ours 的 p95 延迟。在\emph{原执行方式}协议下，每个基线保留其论文中的执行方式，例如在重新规划前执行更长的动作块。

\paragraph{附加延迟。}在推理之后、协议本身的延迟之上，插入 0、50、100、150 或 200\,ms 的附加延迟，期间物理继续推进。

\paragraph{延迟测量。}
\label{app:timing}
我们在一块 \cfgGPU{}（32\,GB，即我们机器人所用的 GPU）上独占运行、以 batch 1 为所有方法计时。这些测量构成仿真 rollout 所回放的延迟分布，独立于 rollout 的计算硬件。每个方法在三轮中各进行三十次预热调用，再在三段各八帧的观测序列上测量 300 次调用；调用边界处做 CUDA 同步，方法顺序轮换。每个方法使用其发布版本的精度，\ours 使用 FP32。出动作时间包括图像编码与生成第一条指令，在真机上还包括跟踪最新一帧。对 \ours，我们另外测量预测头的耗时，尽管部署时跳过它们。我们区分三个时间量。\emph{出动作时间}是从观测到达到其第一条指令可用的实测墙钟时间，不取整；\emph{交接时刻} $\ell_t=k_t\Delta_a$ 把总延迟向上取整到控制周期；\emph{动作年龄}是一条指令的执行时刻减去采集时刻。在真机上，日志记录采集、可用、预处理、推理开始与结束、进入队列、指令下发与执行等时刻。若一次查询的出动作时间超过 50\,ms 的控制周期，下一个周期只能执行剩余指令或保持，我们称该查询超时。附录~\ref{app:latency_results} 报告实测值。

\subsection{基线与消融对照}
\label{app:baselines}
% Protocol table: Dynamic-14 baseline adaptations from BASE1 and the v3 endpoint records.
% PUMA uses the resolved 40k training configuration and parameter summary in
% BASE1/results/puma_dynamic14_20261101/h100/training-final.zip.
\begin{table}[!htbp]
\centering\footnotesize
\caption{\tablelang{\textbf{Baseline training.} The training budget gives optimizer updates and global batch size.
$H$ is the number of commands returned per query. The parameter counts include trainable parameters only.
DOM ckpt. denotes the released DynamicVLA checkpoint trained on DOM.}{\textbf{基线训练。}训练量以优化器更新次数和全局 batch 大小表示。
$H$ 为每次查询返回的指令数。参数量只计入可训练参数。DOM 检查点指 DynamicVLA 发布的、在 DOM 上训练的检查点。}}
\label{tab:baselines}
\setlength{\tabcolsep}{2.6pt}
\begin{tabular}{@{}llccc@{}}
\toprule
\tablelang{Method}{方法} & \tablelang{Initialization}{初始化} & \tablelang{Updates$\times$batch}{更新数$\times$batch} & $H$ & \tablelang{Params}{参数量} \\
\midrule
ACT \cite{act} & \tablelang{scratch}{从头训练} & $50{,}000\times32$ & 100 & 83.9M \\
DP \cite{diffusionpolicy} & \tablelang{scratch}{从头训练} & $50{,}000\times32$ & 15 & 288.7M \\
$\pi_{0.5}$ \cite{pi05} & \texttt{pi05\_base} & $30{,}000\times256$ & 10 & 3,353.4M \\
$\pi_{0.5}$+SE(3) & \texttt{pi05\_base} & $30{,}000\times256$ & 10 & 3,357.7M \\
DynamicVLA \cite{dynamicvla} & \tablelang{DOM ckpt.}{DOM 检查点} & $50{,}000\times44$ & 20 & 99.1M \\
PUMA \cite{domino} & \tablelang{official}{官方} & $40{,}000\times64$ & 16 & 4,516.5M \\
OpenWAM-$\alpha$ \cite{openwam_alpha} & \tablelang{official}{官方} & $34{,}490\times192$ & 32 & 6,021.2M \\
DynamicWAM \cite{dynamicwam} & \tablelang{official}{官方} & $40{,}000\times128$ & 16 & 988.8M \\
DynamicWAM+SE(3) & \tablelang{official}{官方} & $40{,}000\times128$ & 16 & 989.4M \\
\midrule
\ours & \tablelang{scratch}{从头训练} & $\cfgUpdates\times512$ & \cfgChunk & \cfgTrainableParams \\
\bottomrule
\end{tabular}
\end{table}

\paragraph{共同条件。}表~\ref{tab:baselines} 汇总各基线的适配。每个基线保留其官方架构、动作块与初始化。ACT 与 Diffusion Policy 从头训练，$\pi_{0.5}$ 从 \texttt{pi05\_base} 微调，其余方法从其发布的权重出发。它们都在相同的 11{,}200 条示范上训练，示范中的指令通过同一个笛卡尔控制器映射到各自的动作空间。七个已训练基线端点都只使用主视角 RGB，\ours 也只使用主视角。各方法保留原生缺失视图路径：$\pi_{0.5}$ 将两个腕部槽置零并屏蔽；DreamVLA 使用经过预处理的黑色腕部图像并屏蔽腕部辅助监督；OpenWAM-$\alpha$ 与 DynamicWAM 将缺失腕部区域填黑。每个基线沿用官方更新预算，只训练一次，种子为 0；输入与动作适配在本节及表~\ref{tab:baselines} 中说明。对 DreamVLA，我们同时运行了其原始的 40 轮调度与 20 轮余弦调度，并报告在开发场景上更好的一个。

\paragraph{输入。}PUMA \cite{domino} 接收在渲染帧上计算的历史光流。DreamVLA 接收七帧因果 RGB、末端执行器位置与朝向、手指位置和固定的任务文本，用十步 DDIM 和 1.5 的无分类器引导采样。OpenWAM-$\alpha$ 接收当前 RGB 视图、关节值、夹爪开度与任务文本，联合生成视频与动作，但不解码视频。DynamicWAM 接收其光流历史与运动学描述。

\paragraph{加入 SE(3) 的基线。}$\pi_{0.5}$+SE(3) 与 DynamicWAM+SE(3) 从相同的位姿来源接收附录~\ref{app:tokens} 的八个目标状态 token，带有相同的有效标志与掩码。一个两层 MLP 把这些 token 映射到模型宽度。对 $\pi_{0.5}$，我们把它们附加在前缀中的图像 token 之后，动作专家会读取它们；对 DynamicWAM，它们并入其运动学描述。MLP 的输出层零初始化，训练沿用基础模型的配方。

\paragraph{消融对照。}
\label{app:ablation_defs}
第~\mainref{sec:principles}~节的五个对照各改变表~\ref{tab:hyper} 配置中的一个组件，数据、配方、更新预算、执行器以及训练种子 0、1 与 2 保持不变。\emph{无目标历史}去掉目标状态 token 的槽位 0--6，即过去的位姿及其速度与加速度特征，保留存放当前位姿的槽位 7（附录~\ref{app:tokens}）。\emph{无未来监督}去掉 $\mathcal L_P$、$\mathcal L_F$ 与三个未来外观损失及其预测头，保留对当前目标区域的重建。\emph{无场景流监督}去掉 $\mathcal L_F$ 与流读出，\emph{无专用 SE(3)} 去掉 $\mathcal L_P$ 与位姿头；两者都保留目标状态 token 与全部外观目标。\emph{统一精度预测}以\fillin{冻结视频 VAE 在相同时域上对整帧未来隐变量的预测、其读出方式与损失权重}取代 $\mathcal L_P$ 与 $\mathcal L_F$。由于所有预测头在部署时都被跳过，只有无目标历史对照改变推理时的输入。

\subsection{指标与统计}
\label{app:metrics}
按任务等权的成功率为
\begin{equation}
\mathrm{SR}(m)=\frac{100}{14}\sum_{k=1}^{14}\frac1{S_m n_k}\sum_{s=1}^{S_m}\sum_{i=1}^{n_k}y_{mski},
\end{equation}
其中 $y_{mski}$ 是以种子 $s$ 训练的方法 $m$ 在任务 $k$ 的场景 $i$ 上的二值结果；本文各变体 $S_m=3$，基线 $S_m=1$。\emph{首次接触时的位置误差} $e_p$ 是首次接触时末端到专家在目标上的接触坐标系的距离。两个放置任务开始时物体已被抓住，因此其接触是所持物体与容器的首次接触，接触坐标系是专家放置目标所对应的末端位姿。仅含相位阶段的穿越任务（\texttt{t3\_phase}）的成功不需要任何接触，因此 $e_p$ 在其余 13 个任务上平均；它只在发生接触的回合上报告，没有接触的失败仍计入 SR 的分母。

\paragraph{训练种子与区间。}
\label{app:statistics}
完整模型与五个消融对照使用训练种子 0、1 与 2，每个基线使用种子 0。在开发场景上，同一配置的三个训练种子之间 SR 最多相差 \fillin{7.4} 个百分点，主要来自导线穿越，其成功率在各种子间从 \fillin{0} 到 \fillin{71}\% 不等。我们报告每个训练种子以及跨种子的均值与标准差（附录~\ref{app:ablations}）；对消融，我们还在去掉纯相位导线穿越（\texttt{t3\_phase}）后重复比较。对于配对 95\% 区间，我们对每个场景的各种子取平均，并在每个任务内、跨方法联合地对留出场景做自助法。这些区间刻画给定训练权重下的场景不确定性；区间落在 $\pm2$ 个百分点以内的差异视为持平。
